\subsection{Orthonormal Bases and Gram--Schmidt}

Throughout this subsection, $\mathbb{F}$ is $\mathbb{R}$ or $\mathbb{C}$ and $\mathcal{V}$ is an inner-product space over $\mathbb{F}$.

\begin{definition} \label{def:orthogonal-and-orthonormal-sets} \hlt{Orthogonal and Orthonormal Families}\\
For an index set $I$, a family $(\mathbf{u}_i)_{i\in I}$ in $\mathcal{V}$ is \hlt{orthogonal} if $(\forall\,i,j\in I)[i\neq j\Longrightarrow\innerprod{\mathbf{u}_i,\mathbf{u}_j}=0]$. It is \hlt{orthonormal} if $(\forall\,i,j\in I)\,\innerprod{\mathbf{u}_i,\mathbf{u}_j}=\delta_{ij}$.\\
Thus a finite list $(\mathbf{e}_1,\ldots,\mathbf{e}_m)$ is orthonormal exactly when $(\forall\,j,k\in\{1,\ldots,m\})\,\innerprod{\mathbf{e}_j,\mathbf{e}_k}=\delta_{jk}$.\\
Equivalently, the vectors in the list are unit and mutually orthogonal.\\
If $(\mathbf{u}_i)_{i\in I}$ is orthogonal and $(\forall\,i\in I)\,\mathbf{u}_i\neq\mathbf{0}$, then $\left(\frac{\mathbf{u}_i}{\norm{\mathbf{u}_i}}\right)_{i\in I}$ is orthonormal.
\end{definition}

\begin{proposition} \label{prop:orthonormal-linear-combination-norm} \hlt{Norm of an Orthonormal Linear Combination}\\
If $(\mathbf{e}_1,\ldots,\mathbf{e}_m)$ is orthonormal, then $(\forall\,(a_1,\ldots,a_m)\in\mathbb{F}^m)\left[\norm{\sum\limits_{j=1}^{m}a_j\mathbf{e}_j}^2=\sum\limits_{j=1}^{m}\abs{a_j}^2\right]$.
\begin{proof}
$\norm{\sum\limits_{j=1}^{m}a_j\mathbf{e}_j}^2=\innerprod{\sum\limits_{j=1}^{m}a_j\mathbf{e}_j,\sum\limits_{k=1}^{m}a_k\mathbf{e}_k}=\sum\limits_{j=1}^{m}\sum\limits_{k=1}^{m}a_j\overline{a_k}\innerprod{\mathbf{e}_j,\mathbf{e}_k}=\sum\limits_{j=1}^{m}\abs{a_j}^2$ by sesquilinearity and orthonormality.\\
\hlt{Structural Commentary}: Orthonormality deletes every cross term in the squared-norm expansion.\\
\hlt{Key Insight}: In orthonormal coordinates, length is computed coefficient-by-coefficient.
\end{proof}
\end{proposition}

\begin{corollary} \label{cor:orthonormal-list-linearly-independent} \hlt{Orthonormal Lists Are Linearly Independent}\\
If $(\mathbf{e}_1,\ldots,\mathbf{e}_m)$ is orthonormal in $\mathcal{V}$, then it is linearly independent.
\begin{proof}
Suppose $\sum\limits_{j=1}^{m}a_j\mathbf{e}_j=\mathbf{0}$ for $(a_1,\ldots,a_m)\in\mathbb{F}^m$. Then $0=\norm{\mathbf{0}}^2=\sum\limits_{j=1}^{m}\abs{a_j}^2$ by Proposition \ref{prop:orthonormal-linear-combination-norm}.\\
Each $a_j=0$ by nonnegativity of the summands, so the list is linearly independent.\\
\hlt{Structural Commentary}: The norm identity converts a linear relation into a sum of nonnegative scalars.\\
\hlt{Key Insight}: A zero orthonormal linear combination can only have zero coefficients.
\end{proof}
\end{corollary}

\begin{definition} \label{def:orthonormal-basis} \hlt{Orthonormal Basis}\\
An \hlt{orthonormal basis} of $\mathcal{V}$ is a basis of $\mathcal{V}$ that is an orthonormal list.
\end{definition}

\begin{proposition} \label{prop:right-length-orthonormal-list-is-basis} \hlt{Right-Length Orthonormal Lists Are Bases}\\
If $\dim(\mathcal{V})=n<\infty$ and $(\mathbf{e}_1,\ldots,\mathbf{e}_n)$ is orthonormal, then it is an orthonormal basis of $\mathcal{V}$.
\begin{proof}
The list is linearly independent by Corollary \ref{cor:orthonormal-list-linearly-independent}. It is a basis by Proposition \ref{prop:full-dimension-basis-criteria} because its length is $\dim(\mathcal{V})$.\\
\hlt{Structural Commentary}: Orthonormality supplies independence; the dimension criterion supplies spanning.\\
\hlt{Key Insight}: In finite dimensions, $n$ independent vectors in an $n$-dimensional space already form a basis.
\end{proof}
\end{proposition}

\begin{theorem} \label{thm:gram-schmidt-process} \hlt{Gram--Schmidt Process}\\
Let $m\geq1$ and let $(\mathbf{v}_1,\ldots,\mathbf{v}_m)$ be linearly independent in $\mathcal{V}$. Define $\mathbf{u}_1\defeq\mathbf{v}_1$ and $\mathbf{e}_1\defeq\frac{\mathbf{u}_1}{\norm{\mathbf{u}_1}}$.\\
For $2\leq j\leq m$, define $\mathbf{u}_j\defeq\mathbf{v}_j-\sum\limits_{k=1}^{j-1}\innerprod{\mathbf{v}_j,\mathbf{e}_k}\mathbf{e}_k$.\\
For $2\leq j\leq m$, set $\mathbf{e}_j\defeq\frac{\mathbf{u}_j}{\norm{\mathbf{u}_j}}$. Then:
\begin{enumerate}[label=(\roman*)]
\setlength{\itemsep}{0pt}
\item $(\forall\,j\in\{1,\ldots,m\})[\operatorname{span}(\mathbf{v}_1,\ldots,\mathbf{v}_j)=\operatorname{span}(\mathbf{e}_1,\ldots,\mathbf{e}_j)]$;
\item $(\mathbf{e}_1,\ldots,\mathbf{e}_m)$ is orthonormal.
\end{enumerate}
\begin{proof}
\begin{enumerate}[label=(\roman*)]
\setlength{\itemsep}{0pt}
\item[$(i)$] Put $V_j\defeq\operatorname{span}(\mathbf{v}_1,\ldots,\mathbf{v}_j)$ and $E_j\defeq\operatorname{span}(\mathbf{e}_1,\ldots,\mathbf{e}_j)$. Induct on $j$.\\
Since $\mathbf{v}_1\neq\mathbf{0}$, $\mathbf{e}_1$ is a nonzero scalar multiple of $\mathbf{v}_1$, so $V_1=E_1$. Suppose $V_{j-1}=E_{j-1}$.\\
The correction term in $\mathbf{u}_j$ lies in $E_{j-1}=V_{j-1}$, so $\mathbf{u}_j\in V_j$.\\
If $\mathbf{u}_j=\mathbf{0}$, then $\mathbf{v}_j\in V_{j-1}$, contrary to linear independence; hence $\mathbf{u}_j\neq\mathbf{0}$ and $\mathbf{e}_j\in V_j$, so $E_j\subseteq V_j$.\\
Conversely, $\mathbf{v}_j=\norm{\mathbf{u}_j}\mathbf{e}_j+\sum\limits_{k=1}^{j-1}\innerprod{\mathbf{v}_j,\mathbf{e}_k}\mathbf{e}_k\in E_j$, so $V_j\subseteq E_j$.
\item[$(ii)$] $(\forall\,j\in\{1,\ldots,m\})\left[\norm{\mathbf{e}_j}=\norm{\frac{1}{\norm{\mathbf{u}_j}}\mathbf{u}_j}=1\right]$ by Proposition \ref{prop:derived-norm-basic-properties}(iii) and nonvanishing established in (i). Induct on $j$ to verify mutual orthogonality. If $(\mathbf{e}_1,\ldots,\mathbf{e}_{j-1})$ is orthonormal and $k<j$, then $\innerprod{\mathbf{u}_j,\mathbf{e}_k}=\innerprod{\mathbf{v}_j,\mathbf{e}_k}-\sum\limits_{i=1}^{j-1}\innerprod{\mathbf{v}_j,\mathbf{e}_i}\innerprod{\mathbf{e}_i,\mathbf{e}_k}=0$.\\
Hence $\mathbf{e}_j\perp\mathbf{e}_k$ by first-slot homogeneity, and the reverse relation holds by Proposition \ref{prop:orthogonality-and-zero}(i). Thus the full list is orthonormal.
\end{enumerate}
\hlt{Structural Commentary}: Prefix-span induction establishes well-definedness before orthogonality is verified.\\
\hlt{Key Insight}: Each remainder removes precisely the earlier orthonormal components.
\end{proof}
\end{theorem}

\begin{theorem} \label{thm:gram-schmidt-arbitrary-list} \hlt{Gram--Schmidt for Arbitrary Finite Lists}\\
Let $m\geq1$ and let $(\mathbf{v}_1,\ldots,\mathbf{v}_m)$ be a finite list in $\mathcal{V}$. Set $\mathbf{u}_1\defeq\mathbf{v}_1$. For $2\leq k\leq m$, define $\mathcal{J}_{k-1}\defeq\{j\in\{1,\ldots,k-1\}:\mathbf{u}_j\neq\mathbf{0}\}$.\\
For $2\leq k\leq m$, set $\mathbf{u}_k\defeq\mathbf{v}_k-\sum\limits_{j\in\mathcal{J}_{k-1}}\frac{\innerprod{\mathbf{v}_k,\mathbf{u}_j}}{\innerprod{\mathbf{u}_j,\mathbf{u}_j}}\mathbf{u}_j$. Set $U_k\defeq\operatorname{span}(\mathbf{u}_1,\ldots,\mathbf{u}_k)$, $V_k\defeq\operatorname{span}(\mathbf{v}_1,\ldots,\mathbf{v}_k)$, with $U_0=V_0\defeq\{\mathbf{0}\}$. Then:
\begin{enumerate}[label=(\roman*)]
\setlength{\itemsep}{0pt}
\item $(\mathbf{u}_1,\ldots,\mathbf{u}_m)$ is orthogonal;
\item $(\forall\,k\in\{1,\ldots,m\})[U_k=V_k]$.
\end{enumerate}
\begin{proof}
\begin{enumerate}[label=(\roman*)]
\setlength{\itemsep}{0pt}
\item[$(i)$] Induct on $k$. A one-vector list is orthogonal. Suppose $(\mathbf{u}_1,\ldots,\mathbf{u}_{k-1})$ is orthogonal and fix $i<k$.\\
If $\mathbf{u}_i=\mathbf{0}$, then $\mathbf{u}_k\perp\mathbf{u}_i$ by Proposition \ref{prop:orthogonality-and-zero}(ii). If $\mathbf{u}_i\neq\mathbf{0}$, then $i\in\mathcal{J}_{k-1}$ and $\innerprod{\mathbf{u}_k,\mathbf{u}_i}=\innerprod{\mathbf{v}_k,\mathbf{u}_i}-\sum\limits_{j\in\mathcal{J}_{k-1}}\frac{\innerprod{\mathbf{v}_k,\mathbf{u}_j}}{\innerprod{\mathbf{u}_j,\mathbf{u}_j}}\innerprod{\mathbf{u}_j,\mathbf{u}_i}=0$ by the induction hypothesis: all terms with $j\neq i$ vanish, and the $j=i$ term cancels the first term.\\
The reverse orthogonality relation holds by Proposition \ref{prop:orthogonality-and-zero}(i), so the full list is orthogonal.
\item[$(ii)$] Induct on $k$, starting from $U_0=V_0$. Suppose $U_{k-1}=V_{k-1}$.\\
The correction term in $\mathbf{u}_k$ lies in $U_{k-1}=V_{k-1}$, so $\mathbf{u}_k\in V_k$ and $U_k\subseteq V_k$.\\
The recursion also expresses $\mathbf{v}_k$ as $\mathbf{u}_k$ plus a vector in $U_{k-1}$, so $\mathbf{v}_k\in U_k$ and $V_k\subseteq U_k$. Therefore $U_k=V_k$.
\end{enumerate}
\hlt{Structural Commentary}: Orthogonality and prefix-span preservation follow by separate inductions.\\
\hlt{Key Insight}: Subtracting all nonzero earlier directions preserves the generated span.
\end{proof}
\end{theorem}

\begin{corollary} \label{cor:gram-schmidt-zero-remainders} \hlt{Zero Remainders Detect Dependence}\\
For the construction in Theorem \ref{thm:gram-schmidt-arbitrary-list}, $(\forall\,k\in\{1,\ldots,m\})[\mathbf{u}_k=\mathbf{0}\Longleftrightarrow\mathbf{v}_k\in V_{k-1}]$.
\begin{proof}
Fix $k\in\{1,\ldots,m\}$.
\begin{enumerate}[label=\roman*.]
\setlength{\itemsep}{0pt}
\item[$(\Rightarrow)$] Suppose $\mathbf{u}_k=\mathbf{0}$. If $k=1$, then $\mathbf{v}_1=\mathbf{0}\in V_0$. If $k\geq2$, the recursion places $\mathbf{v}_k$ in $U_{k-1}=V_{k-1}$ by Theorem \ref{thm:gram-schmidt-arbitrary-list}(ii).
\item[$(\Leftarrow)$] Suppose $\mathbf{v}_k\in V_{k-1}$. If $k=1$, then $\mathbf{v}_1=\mathbf{0}=\mathbf{u}_1$.\\
If $k\geq2$, then $V_{k-1}=U_{k-1}$ by Theorem \ref{thm:gram-schmidt-arbitrary-list}(ii), so write $\mathbf{v}_k=\sum\limits_{j\in\mathcal{J}_{k-1}}a_j\mathbf{u}_j$ for a coefficient family $(a_j)_{j\in\mathcal{J}_{k-1}}$.\\
The nonzero vectors in this sum are orthogonal by Theorem \ref{thm:gram-schmidt-arbitrary-list}(i), and $(\forall\,i\in\mathcal{J}_{k-1})\,\innerprod{\mathbf{v}_k,\mathbf{u}_i}=a_i\innerprod{\mathbf{u}_i,\mathbf{u}_i}$ by orthogonality.\\
Thus $(\forall\,i\in\mathcal{J}_{k-1})\,a_i=\frac{\innerprod{\mathbf{v}_k,\mathbf{u}_i}}{\innerprod{\mathbf{u}_i,\mathbf{u}_i}}$, so the correction term equals $\mathbf{v}_k$ and $\mathbf{u}_k=\mathbf{0}$ by the recursion.\par
\end{enumerate}
\hlt{Structural Commentary}: Each direction compares the recursion with the preserved prefix span.\\
\hlt{Key Insight}: A zero remainder means the new vector adds no direction.
\end{proof}
\end{corollary}

\begin{figure}[H]
  \centering
  \begin{tikzpicture}[
  x=0.88cm,
  y=0.88cm,
  >=Stealth,
  every node/.style={font=\small},
  subspace/.style={gray!60, line width=0.6pt},
  vec/.style={->, blue!75!black, line width=0.95pt},
  proj/.style={->, teal!70!black, line width=0.85pt},
  residual/.style={->, red!70!black, line width=0.85pt},
  guide/.style={gray!45, dashed, line width=0.45pt},
  point/.style={circle, fill=black, inner sep=1pt}
]
  \coordinate (O) at (0,0);
  \coordinate (P) at (2.55,0);
  \coordinate (V) at (2.55,1.75);
  \coordinate (E) at (2.55,0.95);

  \draw[subspace] (-0.35,0) -- (4.2,0) node[right] {$\operatorname{span}(\mathbf{e}_1,\ldots,\mathbf{e}_{j-1})$};
  \draw[vec] (O) -- (V) node[midway, above left] {$\mathbf{v}_j$};
  \draw[proj] (O) -- (P) node[midway, below] {$\sum_{k<j}\innerprod{\mathbf{v}_j,\mathbf{e}_k}\mathbf{e}_k$};
  \draw[residual] (P) -- (V) node[midway, right] {$\mathbf{u}_j$};
  \draw[residual, red!45!black] (P) -- (E) node[midway, left] {$\mathbf{e}_j$};
  \draw[guide] (V) -- (P);
  \draw[gray!60] (2.35,0) -- (2.35,0.2) -- (2.55,0.2);

  \node[point, label=below left:{$\mathbf{0}$}] at (O) {};
  \node[point] at (P) {};
  \node[point] at (V) {};
\end{tikzpicture}

  \caption{One Gram--Schmidt step: subtract the previous-span component, then normalize.}
  \label{fig:gram-schmidt-step}
\end{figure}

Figure \ref{fig:gram-schmidt-step} depicts the residual after removal of the previous-span component; normalization changes its length while preserving its generated one-dimensional subspace.

\begin{example} \label{ex:gram-schmidt-legendre-polynomials} \hlt{Legendre Polynomials from Gram--Schmidt}\\
Equip $\mathbb{R}[x]$ with $\innerprod{p,q}\defeq\int_{-1}^{1}p(x)q(x)\,dx$ and perform Gram--Schmidt on $(1,x,x^2,x^3,\ldots)$.\\
Parity annihilates mixed even--odd inner products. Since $\innerprod{1,1}=2$, $\innerprod{x^2,1}=\frac{2}{3}$, $\innerprod{x,x}=\frac{2}{3}$, and $\innerprod{x^3,x}=\frac{2}{5}$, the first remainders are $u_1(x)=1$, $u_2(x)=x$, $u_3(x)=x^2-\frac{1}{3}$, and $u_4(x)=x^3-\frac{3}{5}x$.\\
These are nonzero scalar multiples of the first four Legendre polynomials.\\
\hlt{Structural Commentary}: Orthogonal polynomials are orthogonalized monomials.\\
\hlt{Key Insight}: Changing the inner product changes the polynomial family produced by Gram--Schmidt.
\end{example}
\par

\begin{lemma} \label{lem:gram-schmidt-preserves-orthogonal-lists} \hlt{Gram--Schmidt on Orthogonal Lists}\\
Let $m\geq1$ and let $(\mathbf{v}_1,\ldots,\mathbf{v}_m)$ be a mutually orthogonal list of nonzero vectors. Then:
\begin{enumerate}[label=(\roman*)]
\setlength{\itemsep}{0pt}
\item $(\forall\,j\in\{1,\ldots,m\})\,\mathbf{e}_j=\frac{\mathbf{v}_j}{\norm{\mathbf{v}_j}}$;
\item if the original list is orthonormal, then $(\forall\,j\in\{1,\ldots,m\})\,\mathbf{e}_j=\mathbf{v}_j$.
\end{enumerate}
\begin{proof}
\begin{enumerate}[label=(\roman*)]
\setlength{\itemsep}{0pt}
\item[$(i)$] The list is linearly independent by Proposition \ref{prop:mutually-orthogonal-nonzero-vectors-independent}, so the construction is defined by Theorem \ref{thm:gram-schmidt-process}. Induct on $j$.\\
The claim holds for $j=1$. Suppose $(\forall\,i<j)\,\mathbf{e}_i=\frac{\mathbf{v}_i}{\norm{\mathbf{v}_i}}$. Then $(\forall\,i<j)\,\innerprod{\mathbf{v}_j,\mathbf{e}_i}=\frac{1}{\norm{\mathbf{v}_i}}\innerprod{\mathbf{v}_j,\mathbf{v}_i}=0$ by mutual orthogonality.\\
Thus $\mathbf{u}_j=\mathbf{v}_j$ in the recursion and $\mathbf{e}_j=\frac{\mathbf{v}_j}{\norm{\mathbf{v}_j}}$.
\item[$(ii)$] If the input is orthonormal, then $(\forall\,j\in\{1,\ldots,m\})\,\norm{\mathbf{v}_j}=1$ by orthonormality, so $\mathbf{e}_j=\mathbf{v}_j$ by (i).
\end{enumerate}
\hlt{Structural Commentary}: Induction verifies that every orthogonal correction term vanishes.\\
\hlt{Key Insight}: Orthogonal input already has no earlier-direction component.
\end{proof}
\end{lemma}

\begin{corollary} \label{cor:finite-dimensional-inner-product-space-has-orthonormal-basis} \hlt{Existence of Orthonormal Bases}\\
If $\dim(\mathcal{V})<\infty$, then $\mathcal{V}$ has an orthonormal basis.
\begin{proof}
If $\mathcal{V}=\{\mathbf{0}\}$, the empty list is an orthonormal basis. Suppose $\dim(\mathcal{V})=n\geq1$ and choose a basis $(\mathbf{v}_1,\ldots,\mathbf{v}_n)$.\\
By Theorem \ref{thm:gram-schmidt-process}(i)--(ii), its Gram--Schmidt output is orthonormal and has span $\mathcal{V}$, so it is an orthonormal basis.\\
\hlt{Structural Commentary}: Gram--Schmidt orthonormalizes a basis without changing its span.\\
\hlt{Key Insight}: Finite-dimensional inner-product spaces always admit geometry-friendly coordinates.
\end{proof}
\end{corollary}

\begin{corollary} \label{cor:orthonormal-list-extends-to-basis} \hlt{Extending an Orthonormal List}\\
If $\dim(\mathcal{V})<\infty$ and $(\mathbf{e}_1,\ldots,\mathbf{e}_m)$ is orthonormal, then it extends to an orthonormal basis of $\mathcal{V}$.
\begin{proof}
If $m=0$, an orthonormal basis exists by Corollary \ref{cor:finite-dimensional-inner-product-space-has-orthonormal-basis} and extends the empty list.\\
Suppose $m\geq1$. The list is linearly independent by Corollary \ref{cor:orthonormal-list-linearly-independent}. $\exists$ $\mathbf{v}_{m+1},\ldots,\mathbf{v}_n$ such that $(\mathbf{e}_1,\ldots,\mathbf{e}_m,\mathbf{v}_{m+1},\ldots,\mathbf{v}_n)$ is a basis by Theorem \ref{thm:extending-independent-lists}. Let $(\mathbf{f}_1,\ldots,\mathbf{f}_n)$ be its Gram--Schmidt output.\\
The first $m$ steps depend only on the orthonormal prefix, so $(\forall\,j\in\{1,\ldots,m\})\,\mathbf{f}_j=\mathbf{e}_j$ by Lemma \ref{lem:gram-schmidt-preserves-orthogonal-lists}(ii). The output is orthonormal by Theorem \ref{thm:gram-schmidt-process}(ii), and it spans $\mathcal{V}$ by Theorem \ref{thm:gram-schmidt-process}(i) with $j=n$. Thus $(\mathbf{f}_1,\ldots,\mathbf{f}_n)$ is the required extension.\\
\hlt{Structural Commentary}: Extend linearly first, then orthonormalize while preserving the original prefix.\\
\hlt{Key Insight}: Gram--Schmidt leaves an orthonormal prefix fixed.
\end{proof}
\end{corollary}
